\documentclass[11pt]{article}
\usepackage[margin=1in]{geometry}
\usepackage{amsmath,amssymb}
\usepackage{booktabs}
\usepackage{multirow}
\usepackage{hyperref}
\usepackage{xcolor}
\usepackage{caption}
\usepackage{enumitem}
\usepackage{natbib}
\usepackage{graphicx}

\setlength{\parskip}{0.4em}
\setlength{\parindent}{0em}

\title{When Does Symbolic Consensus Help?\\A Systematic Empirical Study}
\author{Anonymous}
\date{}

\begin{document}
\maketitle

\begin{abstract}
Recent proposals augment LLM (Large Language Model) reasoning with symbolic methods that extract formal constraints from reasoning traces and use satisfiability solvers for answer selection.
We present the first systematic empirical evaluation of symbolic consensus, covering three reasoning domains, five model families, and multiple operating conditions.
We instantiate symbolic consensus as SICA (Symbolic Invariant Consensus Aggregation), which extracts first-order logic constraints from $K$ reasoning traces and selects answers via weighted MAX-SAT.
Across more than 1{,}300 problems, SICA provides no significant accuracy gain over standard majority voting (aggregate $|\Delta|\approx1.5$ percentage points).
Extracted constraints are weakly discriminative (mean precision 0.57 vs.\ 0.33 random baseline for 3-class problems), but this modest signal is insufficient to reliably influence answer selection at practical sample sizes.
MAX-SAT optimization is statistically equivalent to simple constraint counting or random weighting.
We identify one boundary condition: at temperature $T{=}1.0$, SICA achieves $+8.3$pp over majority voting ($p{=}0.034$), but extraction quality degrades from 89\% to 76\% and output coherence drops to 85\%, limiting practical value.
We propose $C_0$, the fraction of problems with multiple candidate answers, as a zero-cost binary diagnostic: $C_0{<}5\%$ guarantees method degeneration to majority voting, while $C_0{>}20\%$ is necessary but insufficient for potential gains.
\end{abstract}

%% ===================================================================
\section{Introduction}
%% ===================================================================

Sampling multiple reasoning traces from large language models and aggregating their answers has become standard practice for improving accuracy on complex tasks.
Self-Consistency \citep{wang2023selfconsistency} selects the majority answer across $K$ sampled chains of thought, providing a simple and effective baseline.
Recent work goes further, proposing to extract symbolic constraints from reasoning traces and use formal solvers to aggregate them, potentially exploiting logical structure that majority voting ignores \citep{pan2023logiclm,olausson2023linc}.

The appeal is intuitive: if reasoning traces encode intermediate logical claims, a satisfiability solver should identify the answer most consistent with the full body of extracted knowledge, outperforming a raw head count.
Several methods pursue this direction with varying formalisations of constraint extraction and solver integration.
Yet a basic practical question remains open: across realistic model families, reasoning domains, and operating conditions, when does symbolic consensus actually outperform majority voting?

No prior work systematically characterises the conditions under which symbolic constraint aggregation helps, degenerates to voting, or actively hurts.
Existing evaluations compare symbolic and voting methods on narrow model-domain slices, making it difficult to draw general conclusions about method viability.
This gap is consequential for practitioners who must decide whether the substantial overhead of constraint extraction and solving (approximately 49\% additional compute in our measurements) is justified by accuracy gains.

We address this gap with a controlled empirical study.
We instantiate symbolic consensus as SICA (Symbolic Invariant Consensus Aggregation), which extracts first-order logic (FOL) constraints from $K$ reasoning traces, deduplicates them, and selects answers via weighted MAX-SAT.
SICA serves as a concrete, reproducible test bed for the broader question of whether symbolic aggregation adds value over voting.
We evaluate SICA against majority voting across three reasoning domains (formal logic, commonsense, science), five model families (four instruction-scale, one full reasoning model), and varied operating conditions (temperature $T$, sample count $K$), covering more than 1{,}300 problems with controlled pairwise comparisons.

Our contributions are as follows:
\begin{itemize}[nosep,leftmargin=1.5em]
\item The first systematic evaluation of symbolic consensus across multiple models, domains, and operating conditions, establishing that SICA provides no significant gain over majority voting in standard settings (Section~\ref{sec:main}).
\item Characterisation of the \emph{capability trap} pattern: instruction-scale models produce ample output diversity ($C_0{=}22$--$65\%$), but extracted constraints have only weak discriminative power (precision 0.57), and MAX-SAT over such constraints is equivalent to counting (Sections~\ref{sec:quality}--\ref{sec:ablation}).
\item Analysis of boundary conditions showing that extreme temperature ($T{=}1.0$) partially breaks the pattern but with degraded extraction quality and output coherence (Section~\ref{sec:temperature}).
\item $C_0$, a zero-cost binary diagnostic indicating whether symbolic consensus can potentially operate ($C_0{>}20\%$) or trivially degenerates to voting ($C_0{<}5\%$) (Section~\ref{sec:c0}).
\end{itemize}

%% ===================================================================
\section{Related Work}
%% ===================================================================

\paragraph{Self-Consistency and majority voting.}
Self-Consistency (SC) samples multiple chains of thought and selects the answer receiving the most votes \citep{wang2023selfconsistency}.
Universal Self-Consistency extends this to free-form answers \citep{chen2023universal}.
SC remains the dominant aggregation baseline due to its simplicity and strong empirical performance.
Weighted variants assign non-uniform weights to traces based on length or confidence heuristics \citep{taubenfeld2024weighted}, but evidence for consistent improvement is limited.

\paragraph{Symbolic reasoning for LLMs.}
Logic-LM translates natural language problems into first-order logic and verifies solutions with theorem provers \citep{pan2023logiclm}.
LINC performs neurosymbolic inference over FOL representations \citep{olausson2023linc}.
These methods focus on single-trace symbolic verification rather than multi-trace constraint aggregation.
Our work asks whether the multi-trace aggregation dimension adds value beyond what single-trace verification provides.

\paragraph{Process and outcome reward models.}
Trained verifiers score individual reasoning steps or final answers \citep{cobbe2021gsm8k,lightman2023process}.
Unlike constraint-based aggregation, these methods require supervised training data.
They represent a distinct and complementary approach to improving reasoning quality.

\paragraph{Negative results in NLP.}
Systematic negative-result studies have clarified the limitations of prompting strategies \citep{reynolds2021prompt}, chain-of-thought elicitation, and various reasoning augmentation techniques.
Our work contributes to this tradition by providing controlled evidence on the limits of symbolic consensus.

%% ===================================================================
\section{Method: SICA Pipeline}
\label{sec:method}
%% ===================================================================

SICA (Symbolic Invariant Consensus Aggregation) aggregates multiple reasoning traces through five stages.

\paragraph{1.\ Trace generation.}
Given a problem $q$, we sample $K$ reasoning traces $\{r_1, \ldots, r_K\}$ from an LLM at temperature $T$.
Each trace contains a chain of thought concluding with a candidate answer.

\paragraph{2.\ Answer extraction.}
A deterministic parser extracts the final answer $a_i$ from each trace $r_i$.
Problems where all $K$ traces produce the same answer are classified as \emph{unanimous} and passed directly to majority voting, since no aggregation method can change the outcome.

\paragraph{3.\ Constraint extraction.}
For each trace $r_i$, an LLM-based extractor produces a set of FOL constraints $\mathcal{C}_i = \{c_{i,1}, \ldots, c_{i,m_i}\}$ expressing intermediate logical claims.
Each constraint is a first-order formula over the candidate answer space.
Extraction success is measured as the fraction of traces yielding at least one parseable constraint.

\paragraph{4.\ Deduplication.}
Semantically equivalent constraints across traces are merged, retaining the maximum weight.
Deduplication removes redundancy without discarding signal from independent traces that happen to derive the same intermediate claim.

\paragraph{5.\ MAX-SAT aggregation.}
We construct a weighted MAX-SAT instance: each unique constraint $c_j$ has weight $w_j$ proportional to extraction confidence.
Candidate answers are evaluated by the total weight of constraints they satisfy.
The answer maximising satisfied weight is selected.
When MAX-SAT produces no discrimination (all candidates satisfy the same constraints), the method falls back to majority vote.

\paragraph{$C_0$ metric.}
We define $C_0$ as the fraction of problems in a dataset where $K$ traces produce more than one distinct candidate answer.
$C_0$ reflects output diversity: when $C_0 \approx 0$, all traces agree and any aggregation method reduces to the unanimous answer.
$C_0$ is computable with zero overhead from the answer extraction stage and requires no constraint extraction.

%% ===================================================================
\section{Experimental Setup}
\label{sec:setup}
%% ===================================================================

\paragraph{Models.}
We evaluate seven models spanning four categories.
\emph{Qwen family (instruction-scale)}: Qwen2.5-7B-Instruct, Qwen2.5-14B-Instruct, Qwen2.5-72B-Instruct-AWQ (4-bit quantised).
\emph{Other instruction-scale}: Mistral-7B-Instruct-v0.3, LLaMA-3.1-8B-Instruct.
\emph{Distilled reasoning}: DeepSeek-R1-Distill-Qwen-7B, which we initially expected to behave as a reasoning model but reclassified as instruction-scale after observing $C_0{=}58\%$ (far above the $<5\%$ expected for reasoning models).
\emph{Full reasoning}: QwQ-32B.

\paragraph{Domains.}
ProofWriter (deductive reasoning over formal rules), FOLIO (first-order logic NLI), StrategyQA (commonsense binary reasoning), and ARC-Challenge (science multiple-choice).
The primary evaluation set combines ProofWriter and FOLIO (``D012'') for formal logic coverage.

\paragraph{Default configuration.}
$K{=}12$ traces per problem, temperature $T{=}0.7$, unless stated otherwise.
Temperature and $K$ are varied in ablation experiments.

\paragraph{Statistical testing.}
We use McNemar's exact test for paired comparisons on the same problem set.
For experiments involving multiple comparisons across model families, we apply Bonferroni correction (adjusted $\alpha \approx 0.01$ for five families).

%% ===================================================================
\section{Results}
\label{sec:results}
%% ===================================================================

\subsection{Main Results}
\label{sec:main}

Table~\ref{tab:main} presents SICA versus SC accuracy across all model-domain configurations.

\begin{table}[t]
\centering
\caption{Main results: SICA vs.\ Self-Consistency (SC) across models and domains. $C_0$: fraction of problems with multiple candidate answers. $\Delta$: SICA $-$ SC in percentage points. $p$: McNemar's exact test. Best per-row result in \textbf{bold}.}
\label{tab:main}
\small
\begin{tabular}{@{}llrrrrrr@{}}
\toprule
Model & Domain & $N$ & $C_0$\% & SC & SICA & $\Delta$pp & $p$ \\
\midrule
Qwen-14B & PW+FOLIO & 804 & 38 & \textbf{76.2} & 76.5 & $+0.3$ & n.s. \\
Qwen-14B & StrategyQA & 60 & 10 & 53.3 & \textbf{56.7} & $+3.3$ & 0.16 \\
Qwen-14B & ARC-Chal. & 60 & 5 & \textbf{96.7} & 95.0 & $-1.7$ & --- \\
\midrule
Qwen-7B & PW+FOLIO & 60 & 65 & 56.7 & \textbf{60.0} & $+3.3$ & 0.41 \\
Qwen-72B & PW+FOLIO & 60 & 22 & 85.0 & \textbf{86.7} & $+1.7$ & 0.65 \\
\midrule
Mistral-7B & PW+FOLIO & 60 & 50 & 48.3 & 48.3 & $0.0$ & 1.00 \\
LLaMA-8B & PW+FOLIO & 60 & 62 & 60.0 & \textbf{68.3} & $+8.3$ & 0.096 \\
DeepSeek-R1-7B & PW+FOLIO & 60 & 58 & \textbf{71.7} & 68.3 & $-3.3$ & 0.48 \\
\midrule
QwQ-32B & Mixed & 46 & 2.6 & \multicolumn{2}{c}{$\equiv$ SC} & $0.0$ & --- \\
\bottomrule
\end{tabular}
\end{table}

Across all standard conditions, SICA accuracy falls within a few percentage points of SC.
The largest positive delta ($+8.3$pp, LLaMA-8B) does not reach significance after Bonferroni correction for multiple model comparisons (McNemar $p{=}0.096$, corrected threshold $\approx 0.01$).
This result is further confounded by 48 truncated traces out of 720 total (6.7\% truncation rate), which may artificially inflate diversity.
The largest negative delta ($-3.3$pp, DeepSeek-R1-7B) occurs with the highest truncation rate in the study (72/720 = 10.0\%).
No model-domain combination shows SICA significantly outperforming SC at $\alpha{=}0.05$ after correction.

QwQ-32B represents the converged regime.
With effective $C_0{=}2.6\%$, nearly all problems produce unanimous traces, and SICA trivially reduces to SC.
Full reasoning models, by design, produce deterministic-like chains that eliminate the diversity symbolic methods require.

Within the Qwen family, $C_0$ decreases monotonically with model scale (7B: 65\%, 14B: 38\%, 72B: 22\%, QwQ-32B: 2.6\%), but the SICA-SC gap does not follow a corresponding trend ($+3.3$, $+0.3$, $+1.7$, $0.0$pp).
Diversity reduction does not predict method performance.

DeepSeek-R1-Distill-Qwen-7B was initially expected to function as a reasoning model ($C_0{<}5\%$).
In practice, at 7B parameter capacity, distillation does not produce the deterministic reasoning chains characteristic of full reasoning models.
$C_0{=}58\%$ places it squarely in the instruction-scale cluster, and we reclassify it accordingly.

\subsection{Constraint Quality Analysis}
\label{sec:quality}

Table~\ref{tab:quality} characterises per-problem constraint precision across the 307 multi-candidate problems from the main Qwen-14B experiment ($n{=}804$).

\begin{table}[t]
\centering
\caption{Constraint precision on multi-candidate problems (Qwen-14B, PW+FOLIO, $n{=}307$). Precision measures the fraction of extracted constraints that correctly discriminate the ground-truth answer from alternatives. Random baseline: 0.33 for 3-class, 0.50 for binary.}
\label{tab:quality}
\small
\begin{tabular}{@{}lcccc@{}}
\toprule
Domain & $N$ & Mean & Std & Median \\
\midrule
ProofWriter & --- & 0.55 & --- & --- \\
FOLIO & --- & 0.62 & --- & --- \\
\midrule
Overall & 307 & 0.57 & 0.23 & 0.58 \\
\bottomrule
\end{tabular}
\end{table}

Mean constraint precision is 0.57 (std 0.23), significantly above the 0.33 random baseline for 3-class problems.
Constraints are \emph{weakly discriminative}, not random.
However, precision shows no correlation with vote entropy ($R^2{=}0.008$, $p{=}0.12$), meaning constraint quality does not improve when output diversity increases.
FOLIO constraints (precision 0.62) are somewhat more accurate than ProofWriter constraints (0.55), consistent with FOLIO's more explicit logical structure.

A precision of 0.57 explains the aggregate null result quantitatively.
At $K{=}12$ with $C_0{\approx}38\%$, approximately 4 multi-candidate problems per 60-question set have constraints that could influence the outcome.
With constraints correct 57\% of the time, the expected number of correct answer flips per experiment is approximately 2, which is indistinguishable from sampling noise at these sample sizes.

\subsection{Mechanism Ablation}
\label{sec:ablation}

Table~\ref{tab:ablation} compares four aggregation strategies on the main experiment data ($n{=}804$) to isolate the contribution of MAX-SAT reasoning.

\begin{table}[t]
\centering
\caption{Constraint ablation (Qwen-14B, PW+FOLIO, $n{=}804$, $T{=}0.7$). MAX-SAT: full SICA pipeline. Count-only: constraints weighted equally, no satisfiability reasoning. Random-weight: constraints with random weights. SC: standard majority voting.}
\label{tab:ablation}
\small
\begin{tabular}{@{}lccr@{}}
\toprule
Method & Accuracy (\%) & $\Delta$ vs.\ SC (pp) & $p$ vs.\ SC \\
\midrule
MAX-SAT (SICA) & 76.5 & $+0.3$ & --- \\
Count-only & 76.4 & $+0.2$ & --- \\
Random-weight & 76.3 & $+0.1$ & --- \\
SC (baseline) & 76.2 & --- & --- \\
\midrule
\multicolumn{4}{@{}l}{\small MAX-SAT vs.\ Count-only: $p{=}0.85$; agreement on 802/804 problems.} \\
\multicolumn{4}{@{}l}{\small MAX-SAT vs.\ Random: $p{=}0.72$.} \\
\bottomrule
\end{tabular}
\end{table}

MAX-SAT (76.5\%) is statistically indistinguishable from count-only aggregation (76.4\%), random-weight aggregation (76.3\%), and standard majority voting (76.2\%).
The four methods disagree on only 2 of 804 problems.
The logical structure encoded in constraints is not exploited by the solver in any meaningful way: MAX-SAT degenerates to a constraint counter, and constraints provide no incremental signal beyond raw vote counts.

\subsection{Tiebreak Analysis}
\label{sec:tiebreak}

Table~\ref{tab:tiebreak} examines SICA's performance specifically on tied questions, where majority voting must break ties arbitrarily.

\begin{table}[t]
\centering
\caption{Tiebreak accuracy on tied questions (Qwen-14B, $K{=}12$, $n_{\text{tied}}{=}5$ of 60). Tie fraction: 8.3\%.}
\label{tab:tiebreak}
\small
\begin{tabular}{@{}lc@{}}
\toprule
Tiebreak Method & Accuracy on Ties \\
\midrule
SICA (constraint-based) & 80\% (4/5) \\
Alphabetical & 60\% (3/5) \\
Random (expected) & 50\% \\
First-seen & 40\% (2/5) \\
\bottomrule
\end{tabular}
\end{table}

On the small subset of tied questions (5 of 60, tie fraction 8.3\%), SICA selects the correct answer 80\% of the time versus 50\% for random tiebreaking.
Constraint-based tiebreaking has genuine value on individual ties.
However, ties are too rare to materially affect overall accuracy.
The effective contribution is $0.80 \times 0.083 \approx 6.6\%$ of possible gain, translating to at most ${\sim}2.5$pp under ideal conditions.
At even $K$ (where ties can occur), the tie rate varies from 5.0\% ($K{=}20$) to 26.7\% ($K{=}12$), but even at peak tie rate the overall accuracy impact remains within noise.

\subsection{Boundary Conditions: Temperature}
\label{sec:temperature}

Table~\ref{tab:temperature} presents the temperature ablation on Qwen-14B over 60 D012 problems.
This is the critical experiment for understanding when symbolic consensus can differ from voting.

\begin{table}[t]
\centering
\caption{Temperature ablation (Qwen-14B, D012, $n{=}60$, $K{=}12$). Default operating condition is $T{=}0.7$. Extraction rate: fraction of traces yielding parseable constraints. Coherence: fraction of traces producing coherent, complete outputs.}
\label{tab:temperature}
\small
\begin{tabular}{@{}cccccccc@{}}
\toprule
$T$ & $C_0$\% & SC & SICA & $\Delta$pp & $p$ & Extr. & Coher. \\
\midrule
0.3 & 18 & 73.3 & 75.0 & $+1.7$ & 0.65 & 93\% & ${\sim}98\%$ \\
0.5 & 30 & 73.3 & 75.0 & $+1.7$ & 0.55 & 91\% & ${\sim}98\%$ \\
0.7 & 38 & 76.2 & 76.5 & $+0.3$ & n.s. & 89\% & ${\sim}98\%$ \\
1.0 & 72 & 73.3 & 81.7 & $+8.3$ & \textbf{0.034} & 76\% & 85\% \\
\bottomrule
\end{tabular}
\end{table}

Temperature $T{=}1.0$ produces the only statistically significant SICA advantage in the entire study: $+8.3$pp (McNemar $p{=}0.034$) with a 7:1 discordant ratio favouring SICA.
$C_0$ increases from 38\% at $T{=}0.7$ to 72\% at $T{=}1.0$, confirming the strong link between temperature and diversity ($r{=}0.98$).
The gap is monotonically increasing with temperature.

This gain comes with substantial costs.
Extraction rate drops from 89\% ($T{=}0.7$) to 76\% ($T{=}1.0$): nearly one quarter of constraint extraction attempts produce empty or unparseable output.
Answer coherence drops to 85\%, meaning 15\% of generated traces are partially incoherent.
When extraction retries are factored in, the estimated overhead at $T{=}1.0$ rises to approximately 65\% (vs.\ 49\% at default temperature), increasing the per-question compute cost by roughly 30\%.
The $T{=}0.7$ reference row reports the main experiment ($n{=}804$) for context; all other rows use $n{=}60$.

The practical implication is that the boundary condition exists in theory but is operationally costly.
A practitioner seeking the marginal $+8.3$pp would need to tolerate incoherent outputs, higher extraction failure rates, and additional compute.
For comparison, simply running more traces at the default temperature ($K{=}16$ or $K{=}20$) achieves similar variance reduction without these quality degradations.

Lower temperatures ($T{=}0.3$, $T{=}0.5$) produce small positive deltas ($+1.7$pp each) that are not significant ($p{>}0.5$), consistent with the default-temperature null result.

\subsection{$C_0$ as Operating Regime Indicator}
\label{sec:c0}

Across all 12 model-domain-temperature conditions, $C_0$ values cluster into two distinct regimes (Table~\ref{tab:c0_summary}).

\begin{table}[t]
\centering
\caption{$C_0$ clusters and corresponding SICA-SC gaps. Low-$C_0$ cluster: reasoning models and high-confidence domains. High-$C_0$ cluster: instruction-scale models on challenging domains.}
\label{tab:c0_summary}
\small
\begin{tabular}{@{}lccc@{}}
\toprule
Cluster & $C_0$ range & Mean gap (pp) & $N$ points \\
\midrule
Low $C_0$ & 0--10\% & $-0.9$ & 3 \\
High $C_0$ & 22--72\% & $+1.8$ & 9 \\
\midrule
\multicolumn{4}{@{}l}{\small Linear regression: $R^2{=}0.08$, slope$=0.09$pp per \%$C_0$, $p{=}0.37$.} \\
\bottomrule
\end{tabular}
\end{table}

The low-$C_0$ cluster ($C_0{=}0$--$10\%$) comprises QwQ-32B ($C_0{=}2.6\%$), ARC-Challenge ($C_0{=}5\%$), and StrategyQA ($C_0{=}10\%$).
Mean gap in this cluster is $-0.9$pp: SICA trivially reduces to voting.
The high-$C_0$ cluster ($C_0{=}22$--$72\%$) comprises instruction-scale models on formal logic domains plus the temperature extreme.
Mean gap is $+1.8$pp but with enormous variance (from $-3.3$ to $+8.3$pp), and no individual comparison is significant except $T{=}1.0$.

Linear regression of gap on $C_0$ yields $R^2{=}0.08$ ($p{=}0.37$): $C_0$ provides no gradient-level predictive power for the SICA-SC gap.
The 10--22\% range contains only one data point ($C_0{=}18\%$ at $T{=}0.3$), leaving a gap between clusters.

$C_0$ therefore functions as a binary indicator.
$C_0{<}5\%$ reliably predicts SICA $\equiv$ SC (method degeneration).
$C_0{>}20\%$ indicates that symbolic consensus can potentially differ from voting, but does not predict whether the difference will be positive, negative, or zero.
This makes $C_0$ useful as a zero-cost go/no-go check before investing in constraint extraction, but not as a quantitative performance predictor.

%% ===================================================================
\section{Discussion}
\label{sec:discussion}
%% ===================================================================

\paragraph{The capability trap as descriptive pattern.}
Across our experiments, we observe a recurring pattern that we label the \emph{capability trap}: conditions that enable symbolic consensus to differ from majority voting (high output diversity, $C_0{>}20\%$) do not ensure that the difference is beneficial.
Instruction-scale models produce abundant diversity ($C_0{=}22$--$65\%$), but extracted constraints at precision 0.57 are insufficient to reliably steer answer selection.
We emphasise that this is a descriptive label for an empirical observation, not a mechanistic explanation.
The decoupling of diversity and constraint quality may stem from the extraction process, the FOL formalism, or inherent limitations of current-generation LLM reasoning.

\paragraph{Practical guidelines.}
For practitioners considering symbolic aggregation, our results suggest a simple decision framework.
First, compute $C_0$ from raw answer extractions at zero cost.
If $C_0{<}5\%$ (strong reasoning models, easy tasks), symbolic methods will trivially degenerate to voting.
If $C_0{>}20\%$ under standard temperatures, symbolic methods may differ from voting but are unlikely to improve accuracy, while incurring approximately 49\% additional compute (Appendix~\ref{app:overhead}).

\paragraph{The $T{=}1.0$ exception and its limits.}
The temperature ablation reveals that extreme diversity ($T{=}1.0$, $C_0{=}72\%$) can produce a significant SICA advantage.
This finding prevents us from claiming that the capability trap is absolute.
However, the associated quality degradation (extraction rate 76\%, coherence 85\%) and the fact that this result comes from a single $n{=}60$ experiment temper its practical significance.
Replication across models and larger sample sizes is needed before recommending $T{=}1.0$ as an operating point.

\paragraph{Alternative aggregation directions.}
The Logic-LM comparison (Appendix~\ref{app:logiclm}) shows that Logic-LM with $K$ repetitions and majority voting achieves 80.9\% on FOLIO vs.\ SICA's 79.4\%.
Single-trace FOL verification repeated $K$ times may be more effective than multi-trace constraint aggregation.
Inverse-length weighted voting (Appendix~\ref{app:weightedsc}) matches SICA at zero extraction cost (77.0\% vs.\ 76.5\%).
These results suggest that lightweight quality proxies, or per-trace verification, may be more productive directions than cross-trace constraint consensus.

\paragraph{Limitations.}
(1) We evaluate one instantiation of symbolic consensus (FOL constraints, MAX-SAT solver). Other constraint languages or aggregation strategies might yield different results.
(2) The largest powered test ($n{=}804$) covers only Qwen-14B on ProofWriter and FOLIO. Cross-model tests at $n{\geq}200$ exist only for LLaMA-8B.
(3) We test only one full reasoning model (QwQ-32B); additional reasoning models (o1-series, Gemini thinking models) would strengthen the converged-regime claim.
(4) The $T{=}1.0$ boundary condition has been tested only on Qwen-14B with $n{=}60$.
(5) The constraint ablation (Table~\ref{tab:ablation}) was performed at $T{=}0.7$; results may differ at extreme temperatures where the constraint quality distribution changes.

%% ===================================================================
\section{Conclusion}
%% ===================================================================

We present the first systematic empirical study of symbolic consensus for LLM reasoning aggregation.
Across three domains, five model families, and over 1{,}300 problems, symbolic constraint aggregation via SICA provides no significant accuracy improvement over majority voting in standard operating conditions.
Extracted constraints are weakly discriminative (precision 0.57) but insufficient to reliably influence answer selection.
The MAX-SAT solver degenerates to a constraint counter: it is statistically equivalent to counting or random weighting.
At extreme temperatures ($T{=}1.0$), a marginal gain appears ($+8.3$pp, $p{=}0.034$) but is accompanied by degraded extraction quality and output coherence.
We contribute $C_0$ as a practical, zero-cost binary diagnostic for practitioners deciding whether to invest in symbolic aggregation.
Our findings indicate that, for the current generation of LLM constraint extraction, simple majority voting remains the method of choice.

%% ===================================================================
%% APPENDIX
%% ===================================================================
\appendix

\section{$K$ Sensitivity Analysis}
\label{app:k_ablation}

Table~\ref{tab:k_ablation} presents the effect of varying $K$ on SICA vs.\ SC (Qwen-14B, D012, $n{=}60$).

\begin{table}[h]
\centering
\caption{$K$ ablation (Qwen-14B, D012, $n{=}60$, $T{=}0.7$). Tiebreak rate: fraction of problems producing exact ties at this $K$.}
\label{tab:k_ablation}
\small
\begin{tabular}{@{}rccccr@{}}
\toprule
$K$ & SICA & SC & $\Delta$pp & Tiebreak Rate \\
\midrule
6  & 75.0 & 73.3 & $+1.7$ & 18.3\% \\
8  & 75.0 & 75.0 & $0.0$  & --- \\
11 & 75.0 & 75.0 & $0.0$  & --- \\
12 & 73.3 & 75.0 & $-1.7$ & 26.7\% \\
16 & 71.7 & 73.3 & $-1.7$ & --- \\
20 & 71.7 & 73.3 & $-1.7$ & 5.0\% \\
\bottomrule
\end{tabular}
\end{table}

At odd $K{=}11$, where exact ties cannot occur, the gap is precisely zero.
This confirms that SICA's potential value comes entirely from tiebreaking.
At $K{=}6$ (small sample), the $+1.7$pp positive delta falls within the bootstrap 95\% CI $[-5.0, +8.3]$pp, compatible with zero effect.
At $K{\geq}12$, SICA shows a consistent $-1.7$pp deficit, likely from noise injection by low-quality constraints in larger trace sets.
The pattern is non-monotonic (positive at low $K$, zero at odd $K$, negative at high $K$), making a blanket ``SICA $\leq$ SC at all $K$'' statement inaccurate.
The correct characterisation is \emph{$K$-dependent with no consistent advantage}.

\section{Overhead Analysis}
\label{app:overhead}

Table~\ref{tab:overhead} breaks down wall-clock time per question from the main experiment ($n{=}804$).

\begin{table}[h]
\centering
\caption{Per-question wall-time breakdown (Qwen-14B, $n{=}804$, $T{=}0.7$).}
\label{tab:overhead}
\small
\begin{tabular}{@{}lrr@{}}
\toprule
Component & Mean time (s) & \% of total \\
\midrule
Trace generation & 31.6 & 51.3\% \\
Constraint extraction & 29.0 & 47.0\% \\
Deduplication & 1.0 & 1.6\% \\
MAX-SAT solving & 0.003 & ${<}0.01\%$ \\
\midrule
Total & 61.6 & 100\% \\
\bottomrule
\end{tabular}
\end{table}

SICA's overhead relative to SC (which requires only trace generation) is 48.7\%, dominated by the constraint extraction LLM call.
MAX-SAT solving time is negligible.
The overhead-to-gain ratio at default temperature is $48.7\% / 0.3\text{pp} \approx 162$: each percentage point of (non-significant) gain costs 162\% in additional compute.
Median overhead is 48.9\%, with the 90th percentile at 57.6\%.
ProofWriter questions incur 51.3\% overhead on average vs.\ 38.6\% for FOLIO, reflecting longer constraint extraction prompts for deductive rules.

\section{Weighted SC Baselines}
\label{app:weightedsc}

Table~\ref{tab:weighted_sc} compares SICA against simple weighted voting heuristics ($n{=}804$).

\begin{table}[h]
\centering
\caption{Weighted SC variants (Qwen-14B, PW+FOLIO, $n{=}804$, $T{=}0.7$).}
\label{tab:weighted_sc}
\small
\begin{tabular}{@{}lccc@{}}
\toprule
Method & Accuracy (\%) & $\Delta$ vs.\ SC (pp) & $p$ vs.\ SC \\
\midrule
Inverse-length & 77.0 & $+0.8$ & 0.42 \\
SICA (MAX-SAT) & 76.5 & $+0.3$ & --- \\
SC (uniform) & 76.2 & --- & --- \\
Length-weighted & 75.3 & $-0.9$ & 0.16 \\
Constraint-count & 75.3 & $-0.9$ & --- \\
\bottomrule
\end{tabular}
\end{table}

Inverse-length weighting (favouring shorter, more confident traces) achieves 77.0\%, slightly above SICA's 76.5\%.
All differences are non-significant.
A zero-cost heuristic based on trace length matches or exceeds the full SICA pipeline, reinforcing the conclusion that constraint extraction does not add value at default settings.

\section{Logic-LM Comparison}
\label{app:logiclm}

Table~\ref{tab:logiclm} compares SICA against a Logic-LM baseline on FOLIO ($n{=}204$).

\begin{table}[h]
\centering
\caption{Logic-LM $\times$ $K$ vs.\ SICA vs.\ SC on FOLIO ($n{=}204$, $K{=}12$).}
\label{tab:logiclm}
\small
\begin{tabular}{@{}lcc@{}}
\toprule
Method & Accuracy (\%) & $p$ vs.\ SICA \\
\midrule
Logic-LM $\times$ $K$ + vote & 80.9 & 0.38 \\
SICA (MAX-SAT) & 79.4 & --- \\
SC (majority vote) & 79.4 & --- \\
Single-trace Logic-LM & 77.5 & --- \\
\bottomrule
\end{tabular}
\end{table}

Logic-LM repeated $K$ times with majority voting achieves 80.9\% vs.\ SICA's 79.4\%.
The difference is not significant ($p{=}0.38$), but the directional advantage of per-trace verification over multi-trace constraint aggregation is notable.
Logic-LM performs a complete FOL consistency check per trace (pass/fail), which has stronger discriminative power than SICA's soft constraint aggregation: an internally inconsistent trace is excluded entirely by Logic-LM, while SICA still incorporates its constraints.

\section{GSM8K Pilot}
\label{app:gsm8k}

A 30-question pilot on GSM8K (Qwen-14B, $K{=}12$) yields: SICA 80.0\% vs.\ SC 76.7\% ($\Delta{=}+3.3$pp, $p{=}0.56$, McNemar).
$C_0{=}27\%$ (moderate), multi-candidate rate 23\%, extraction rate 85\%.
At $n{=}30$, the standard error is approximately 8pp, making this result uninformative.
GSM8K is included as a pilot data point for domain coverage; it does not alter any main-text claims.

\section{Constraint Case Studies}
\label{app:case_studies}

We qualitatively analyse 8 representative questions from the main experiment.
Of these, 3 produce semantically meaningful constraints, 2 produce correctly discriminative constraints that steer SICA to the right answer, 3 produce non-discriminative constraints (tautological or irrelevant), and 2 produce actively misleading constraints that steer SICA away from the correct answer.

Constraint categories: 3 tautological, 2 partially correct, 2 fully correct, 1 contradicting ground truth.
The two fully correct cases both come from FOLIO problems with explicit formal logical structure, where FOL extraction naturally aligns with the problem's own representation.
The misleading case involves a ProofWriter problem where the extracted constraint encodes a premise not present in the original rules.

These case studies illustrate that constraints \emph{can} be discriminative on individual problems (consistent with the aggregate precision of 0.57), but the mixture of helpful, neutral, and harmful constraints yields no net improvement when aggregated across a full problem set.

\section{$C_0$ Diagnostic Full Data}
\label{app:c0_data}

Table~\ref{tab:c0_full} lists all 12 data points used for the $C_0$ analysis in Section~\ref{sec:c0}.

\begin{table}[h]
\centering
\caption{All (model/condition, domain, $C_0$, SICA-SC gap) data points.}
\label{tab:c0_full}
\small
\begin{tabular}{@{}llrr@{}}
\toprule
Condition & Domain & $C_0$ (\%) & Gap (pp) \\
\midrule
QwQ-32B & Mixed & 2.6 & $0.0$ \\
Qwen-14B & ARC & 5.0 & $-1.7$ \\
Qwen-14B & StrategyQA & 10.0 & $+3.3$ \\
Qwen-14B $T{=}0.3$ & PW+FOLIO & 18.0 & $+1.7$ \\
Qwen-72B & PW+FOLIO & 22.0 & $+1.7$ \\
Qwen-14B GSM8K & GSM8K & 27.0 & $+3.3$ \\
Qwen-14B $T{=}0.7$ & PW+FOLIO & 38.0 & $+0.3$ \\
Mistral-7B & PW+FOLIO & 50.0 & $0.0$ \\
DeepSeek-R1-7B & PW+FOLIO & 58.0 & $-3.3$ \\
LLaMA-8B & PW+FOLIO & 62.0 & $+8.3$ \\
Qwen-7B & PW+FOLIO & 65.0 & $+3.3$ \\
Qwen-14B $T{=}1.0$ & PW+FOLIO & 72.0 & $+8.3$ \\
\bottomrule
\end{tabular}
\end{table}

%% ===================================================================
%% REFERENCES
%% ===================================================================

\bibliographystyle{plainnat}
\begin{thebibliography}{20}

\bibitem[Wang et~al.(2023)]{wang2023selfconsistency}
X.~Wang, J.~Wei, D.~Schuurmans, Q.~Le, E.~Chi, S.~Narang, A.~Chowdhery, and D.~Zhou.
\newblock Self-Consistency Improves Chain of Thought Reasoning in Language Models.
\newblock In \emph{ICLR}, 2023.

\bibitem[Chen et~al.(2023)]{chen2023universal}
X.~Chen, R.~Aksitov, U.~Alon, J.~Ren, K.~Xiao, P.~Yin, S.~Prakash, C.~Sutton, X.~Wang, and D.~Zhou.
\newblock Universal Self-Consistency for Large Language Model Generation.
\newblock \emph{arXiv:2311.17311}, 2023.

\bibitem[Pan et~al.(2023)]{pan2023logiclm}
L.~Pan, A.~Albalak, X.~Wang, and W.~Y.~Wang.
\newblock Logic-LM: Empowering Large Language Models with Symbolic Solvers for Faithful Logical Reasoning.
\newblock In \emph{Findings of EMNLP}, 2023.

\bibitem[Olausson et~al.(2023)]{olausson2023linc}
T.~X.~Olausson, A.~Gu, B.~Lipkin, C.~E.~Zhang, A.~Solar-Lezama, J.~B.~Tenenbaum, and R.~Levy.
\newblock LINC: A Neurosymbolic Approach for Logical Reasoning by Combining Language Models with First-Order Logic Provers.
\newblock In \emph{EMNLP}, 2023.

\bibitem[Cobbe et~al.(2021)]{cobbe2021gsm8k}
K.~Cobbe, V.~Kosaraju, M.~Bavarian, M.~Chen, H.~Jun, L.~Kaiser, M.~Plappert, J.~Tworek, J.~Hilton, R.~Nakano, C.~Hesse, and J.~Schulman.
\newblock Training Verifiers to Solve Math Word Problems.
\newblock \emph{arXiv:2110.14168}, 2021.

\bibitem[Lightman et~al.(2023)]{lightman2023process}
H.~Lightman, V.~Kosaraju, Y.~Burda, H.~Edwards, B.~Baker, T.~Lee, J.~Leike, J.~Schulman, I.~Sutskever, and K.~Cobbe.
\newblock Let's Verify Step by Step.
\newblock In \emph{ICLR}, 2024.

\bibitem[Taubenfeld et~al.(2024)]{taubenfeld2024weighted}
A.~Taubenfeld, Y.~Horesh, N.~Benaim, G.~Shani, N.~Shani, and L.~Wolf.
\newblock Systematic Biases in LLM Simulations of Debates.
\newblock \emph{arXiv:2402.04049}, 2024.

\bibitem[Reynolds and McDonell(2021)]{reynolds2021prompt}
L.~Reynolds and K.~McDonell.
\newblock Prompt Programming for Large Language Models: Beyond the Few-Shot Paradigm.
\newblock In \emph{CHI EA}, 2021.

\end{thebibliography}

\end{document}
